% !TeX program = lualatex
% halo:begin
% {
%   "version": 1,
%   "publish": false,
%   "course": "data-and-optimization",
%   "section": "1.1",
%   "title": "数据与优化（通识笔记）1.1：大数据与最优化",
%   "categories": ["study-notes"],
%   "tags": ["mathematics"],
%   "excerpt": "大数据基础与思维、机器学习和深度学习、优化建模、回归与分类、稀疏与低秩结构，以及优化算法与收敛。"
% }
% halo:end
% 来源：通识课-1大数据案例介绍.pdf（1--60页）；通识课-2大数据思维.pdf（1--33页）；
% 01-opt-dzw.pdf（1--45页）；05-lect1-model.pdf（1--31页，支持向量机结束）。
% 课程范围由用户于2026-10-28确认；教材不纳入复习范围。
% 来源定位与写作阶段记录见Beads SHU-notes-0w4。
% 信息论公式、统一符号及简短推导用于解释课件概念，不另扩展教材章节。
\RequirePackage{iftex}
\RequireLuaTeX
\documentclass[UTF8,fontset=fandol,a4paper,12pt]{ctexart}
\usepackage[margin=2.5cm]{geometry}
\usepackage{amsmath,amssymb,amsthm}
\usepackage{enumitem,fancyhdr,booktabs,longtable,array}
\usepackage[hidelinks]{hyperref}
\newcommand{\R}{\mathbb R}
\newcommand{\E}{\mathbb E}
\newcommand{\rank}{\operatorname{rank}}
\newcommand{\vx}{\boldsymbol x}
\newcommand{\va}{\boldsymbol a}
\newcommand{\vb}{\boldsymbol b}
\newcommand{\vw}{\boldsymbol w}
\newcommand{\norm}[1]{\left\lVert#1\right\rVert}
\theoremstyle{definition}
\newtheorem{definition}{定义}[section]
\newtheorem{proposition}[definition]{命题}
\newtheorem{example}[definition]{例}
\renewcommand{\proofname}{证明}
\setlist[enumerate]{leftmargin=2em,itemsep=0.3em,topsep=0.3em}
\setlist[itemize]{leftmargin=2em,itemsep=0.2em,topsep=0.3em}
\setlength{\headheight}{16pt}
\pagestyle{fancy}\fancyhf{}
\fancyhead[L]{\small 数据与优化（通识笔记）}
\fancyhead[R]{\small 大数据与最优化}
\fancyfoot[C]{\thepage}
\renewcommand{\headrulewidth}{0.3pt}
\raggedbottom
\title{大数据与最优化}
\author{Yuchen Zhang}
\date{整理日期：2026年10月28日}
\makeatletter
\renewcommand{\maketitle}{\begin{center}
{\LARGE\bfseries\@title\par}\vspace{0.6em}
\@author\par\vspace{0.2em}{\small\@date\par}
\end{center}}
\makeatother
\begin{document}
\maketitle
评分记录可以用于推荐，搜索记录可以辅助预测，交通记录可以支持路线选择。数据的作用取决于所研究的问题，也取决于怎样从记录中提取信息、建立模型并检验结果。
\tableofcontents
\newpage

\section{数据、信息与大数据}
\subsection{从记录到知识}
\begin{definition}
\textbf{数据}是对事物的状态、属性或变化的记录，可以表现为数值、文字、图像、声音和行为记录等。\textbf{信息}是数据在具体问题中表达的意义。\textbf{知识}是对信息进一步组织、解释后形成的系统认识。
\end{definition}

例如，天体在各个时刻的位置是数据，由这些记录得到的运动轨迹是信息，对运动规律的概括则形成知识。由观测建立模型，再用模型预测新的观测，构成了从数据取得认识的一条基本途径。

同样多的记录，所能提供的信息未必相同。重复一条已知记录，通常不会增加同等数量的新信息；与研究目标无关的数据，也未必能改善判断。因此，数据的价值既取决于记录的内容和质量，也取决于所要解决的问题。

\subsection{大数据的特点与来源}
大数据的困难不仅在于数量多，还在于数据产生快、类型多，并且有用信息需要从大量记录中提取。其基本特点常概括为\textbf{4V}。
\begin{center}
\begin{tabular}{p{0.23\textwidth}p{0.66\textwidth}}
\toprule
特点 & 含义及处理要求\\
\midrule
规模巨大（Volume） & 数据量超出常规处理方式的有效承受范围，需要考虑存储、计算与扩展能力。\\
高速产生（Velocity） & 数据持续生成、更新，分析结果常有时效要求，迟到的结果可能失去作用。\\
形式多样（Variety） & 数值表格、文本、图像、音频、视频等并存，需要不同的表示和处理方法。\\
潜藏价值（Value） & 大量记录中蕴含有用规律，但价值密度可能较低，需要围绕任务进行分析。\\
\bottomrule
\end{tabular}
\end{center}

数据库中的规则表格属于结构化数据，文章、影像和语音通常属于非结构化数据。传感器、移动终端、社交平台和交易系统不断产生这些记录，移动互联网又使不同记录能够相互关联。

\textbf{大数据与单纯的大规模数据有所区别。}除了体量，还要考察数据是否具有多维关联，是否覆盖研究所需的情形，以及能否及时更新。例如，同一条交通记录同时包含位置、时间、速度和行驶方向，就能支持比单独统计车辆总数更丰富的判断。这里的完备性指与任务有关的覆盖程度，并不意味着掌握了所有事实。

\subsection{数据分析的完整过程}
从数据到结论，需要经过问题定义、数据处理、建模和验证。
\begin{enumerate}
\item \textbf{确定问题。}明确要描述什么、预测什么或选择什么，并规定评价标准。
\item \textbf{收集与处理数据。}检查来源、缺失、重复、异常和测量误差，统一记录的时间、单位与口径。
\item \textbf{形成表示。}把原始记录转化为模型能够利用的样本、特征或矩阵。
\item \textbf{选择模型并估计参数。}从较简单的模型起步，用数据判断是否有必要增加复杂性。
\item \textbf{验证与比较。}在相同条件下比较候选方法，检验新数据上的误差、计算代价与适用范围。
\item \textbf{应用与反馈。}把结果用于预测或决策，并在环境变化后重新检查模型。
\end{enumerate}
数据收集的\textbf{量}与\textbf{质}必须同时考虑。错误记录和缺乏代表性的样本，即使数量很多，也可能使后续分析偏离实际目标。

在这一过程中，数据库技术负责组织、存储和查询记录，统计学研究数据中的变异、估计与推断，数据挖掘侧重发现有用模式，机器学习侧重从数据学习预测或表示。这些领域相互交叉；数据挖掘的英文为 data mining，深度学习为 deep learning，两者含义不同。

\section{大数据思维与不确定性}
\subsection{确定性描述与概率描述}
确定性模型用规律描述变量之间的关系；当影响因素复杂或观测不完整时，概率模型用分布表达可能结果。

机械思维强调：规律是确定的，可以用较简单的公式或语言描述，并能在相应条件下用于预测。其优势是结构清楚、解释明确；但当变量很多、关系复杂、观测不充分时，直接写出完整的确定性规律可能十分困难。

大数据思维承认判断中的不确定性，利用大量观测、多维关联和统计规律改善预测。\textbf{数据驱动}是从样本中学习规律或参数；\textbf{机理驱动}是从对象的内在机制建立关系。二者可以结合，例如用物理规律限制模型，再用数据估计未知参数。数据驱动不要求放弃解释，概率描述也不等于现象没有规律。

\subsection{信息熵与信息的作用}
信息的作用可以从不确定性的变化来理解。设离散随机变量 $X$ 的可能取值为 $x_1,\ldots,x_s$，相应概率为 $p_1,\ldots,p_s$。
\begin{definition}[信息熵]
\[
 \boxed{H(X)=-\sum_{i=1}^s p_i\log_2 p_i,}
\]
其中约定 $0\log_2 0=0$。以 $2$ 为底时，单位为比特。
\end{definition}
若结果已经确定，则熵为 $0$；若 $s$ 个结果等可能，则熵为 $\log_2 s$，也是固定取值个数下的最大熵。例如两个等可能结果的熵为 $1$ 比特。熵越大，事先判断结果的不确定性越大。

取得额外观测 $Y$ 后，平均剩余的不确定性用条件熵 $H(X\mid Y)$ 表示。互信息
\[
 I(X;Y)=H(X)-H(X\mid Y)\ge0
\]
表示观测 $Y$ 平均减少了多少关于 $X$ 的不确定性。若 $Y$ 与 $X$ 独立，则互信息为 $0$；若新增记录只是已有信息的重复，通常也不能提供同等数量的新信息。例如，判断道路是否拥堵时，增加车辆位置和速度记录可能有用，反复保存同一时刻的同一位置则主要增加冗余。

\subsection{编码、传输与概率模型}
信息熵不仅解释不确定性，还给出了无损编码的基本界限。对于独立同分布的离散信源，采用可唯一译码的编码时，平均每个符号的码长不能低于信源熵；通过对足够长的符号块联合编码，可以使平均每个符号的码长任意接近这一界限。这是信源编码定理所说明的压缩极限。

信道编码定理讨论有噪声的传输。在相应的信道模型下，低于信道容量的传输速率可以通过适当编码实现任意低的错误率；超过容量则不能保证可靠传输。\textbf{前者讨论怎样压缩，后者讨论怎样可靠传送}，两者所回答的问题不同。

当只掌握某些统计条件时，可以在满足这些条件的概率分布中选择熵最大的一个，这称为\textbf{最大熵原理}。它保留已知信息，同时避免对未知部分施加额外的无依据限制；并不表示可以忽略已有知识。

比较真实分布 $P=(p_i)$ 与模型分布 $Q=(q_i)$ 时，交叉熵与相对熵分别为
\[
 H(P,Q)=-\sum_i p_i\log_2 q_i,\qquad
 D_{\mathrm{KL}}(P\Vert Q)=\sum_i p_i\log_2\frac{p_i}{q_i}.
\]
在 $p_i>0$ 的位置要求 $q_i>0$，否则相应交叉熵为无穷大。两者满足
\[
 \boxed{H(P,Q)=H(P)+D_{\mathrm{KL}}(P\Vert Q).}
\]
因此，$P=Q$ 时相对熵为 $0$，而交叉熵等于 $H(P)$，一般并不为 $0$。交叉熵可以衡量用模型分布解释数据所需的平均编码代价，后文分类模型的对数损失与此相联系。

\subsection{相关关系与因果关系}
相关关系描述共同变化，因果关系则涉及一个因素的改变怎样影响另一个因素。例如，某类搜索量与流感活动可能同时上升，因此搜索量能够成为预测线索；但人为增加搜索次数，并不意味着会增加患病人数。

两个变量相关，可能是一个影响另一个，也可能是反向作用，或二者受共同因素影响。因此，\textbf{相关关系可以支持预测，却不能单独证明因果关系}。要从预测转向干预，还需要相应的研究设计和证据。多来源记录之间的相互核验可以提高信息可信度，但“交叉核验”与后文用于评价模型的“交叉验证”是不同概念。

\subsection{数据量、代表性与误差}
增加样本量可以减小某些随机误差，却不能自动消除采样偏差。设随机变量 $X$ 有有限期望 $\mu$ 和方差 $\sigma^2$，则切比雪夫不等式给出
\[
 \Pr\{|X-\mu|\ge\varepsilon\}\le\frac{\sigma^2}{\varepsilon^2},
 \qquad\varepsilon>0.
\]
若 $X_1,\ldots,X_N$ 独立同分布，具有上述期望和方差，则样本均值 $\overline X_N$ 满足
\[
 \Pr\{|\overline X_N-\mu|\ge\varepsilon\}
 \le\frac{\sigma^2}{N\varepsilon^2}\longrightarrow0.
\]
这里趋于稳定的是\textbf{样本均值}，不是每条观测的噪声都消失了。切比雪夫不等式可以用于证明这一形式的大数定律，但不等式与大数定律本身不是同一个命题。

若样本只来自某类平台用户，就可能不能代表全部人群；若测量设备始终偏高，增加记录也不能消除这种系统偏差。模型还可能因时间、地区或行为习惯变化而失效，这称为分布变化。使用数据时，除了预测误差，还应考虑采集授权、隐私及对不同群体的影响。

\section{大数据应用与定量决策}
\subsection{推荐、搜索与预测}
推荐系统需要从已知偏好推测未知偏好。评分矩阵的行对应用户，列对应物品，元素记录评分；用户信息矩阵记录年龄、职业等属性，物品信息矩阵记录类型、内容等属性。后两者能够提供辅助信息，不同推荐方法对这些信息的依赖不同。

\textbf{评分缺失不等于评分为零。}缺失表示没有观测，零分则是一种已知评价。协同过滤利用不同用户或物品之间的共同结构进行预测，低秩矩阵补全是其中一种建模途径，见第\ref{sec:lowrank}节。预测后再按照偏好、候选范围和应用要求形成推荐。

搜索记录能够提供及时的行为信号，但搜索行为与目标现象之间的关系仍须检验。流感预测案例的基本过程是：选择可能相关的搜索特征，利用历史搜索量与实际疾病观测训练模型，再与新的实际观测比较。其作用是补充较及时的趋势信号；新闻关注和搜索习惯也会改变搜索量，因此不能把搜索次数直接当作患病人数。

广告匹配同样利用用户、内容和点击记录估计相关程度，购物关联分析则寻找共同出现的购买模式。药物用途探索从已有记录中寻找新的候选联系，再交由后续实验验证。这些应用共同体现了\textbf{从大量记录发现线索，再通过验证判断线索能否使用}的过程。

\subsection{多源数据与实时判断}
交通管理需要同时了解不同路段的状态。传统方法依赖固定观察、上报和汇总，覆盖与更新速度受到限制；移动终端产生的位置、速度和时间记录，使动态路况估计成为可能。数据的价值来自覆盖范围、实时性和相互关联，也取决于能否区分真实拥堵与停车、定位误差等其他情况。

饮食问答除了问题文字，还带有地域、时间和用户背景等维度。按这些维度组织记录，可以比较不同地区和群体的关注点；但提问者的关注不能直接等同于总体人群的实际饮食比例。选举等群体预测也要区分观测样本与目标人群，不能把网络发言简单当作全体人的选择。

异常检测则把个案与同类对象的统计规律比较。例如，用电模式异常可以作为进一步调查的线索，企业收入与同规模、同地区企业差异较大也可能需要核查。\textbf{异常不等于违规，统计筛查也不等于事实认定}；数据分析首先帮助缩小核查范围。

多源信息还可以相互补充。大庆油田的案例中，公开照片、新闻报道等分别提供间接线索，把这些线索放在一起，便能推断单条材料未直接说明的情况。这里的关键是比较不同来源能否相互印证，并区分已知事实与据此作出的推断。

\subsection{行为记录、社会关系与隐私}
用户行为不仅反映个人偏好，也可能受所处环境影响。锚定效应指先接触的数值或评价会影响后续判断；推荐中的评分因此未必是彼此独立、完全稳定的偏好测量。把这种行为机制纳入分析，有助于解释评分结构，并检验是否能改善推荐误差。

通话记录还能构成社会关系网络：用户是节点，联系是边，联系频率等可作为边的属性。网络联系的多样性可以与地区经济指标比较，通话模式也可能与年龄、性别或关系类型有关。这些案例把个体行为转成群体结构特征，展示了数据的多维价值；相关图形本身仍不能证明社会关系多样性直接导致经济增长。

跨域预测进一步说明，数据的潜在用途可能超出采集时的表面目的。由用户点赞构成矩阵，再用奇异值分解提取低维特征，能够把这些特征送入逻辑回归或线性回归，预测其他属性。这一过程把\textbf{行为数据、低维表示、监督学习和验证}连接起来，也说明看似普通的公开行为可能泄露未直接填写的私人信息。

评价时还应辨认指标。均方根误差
\[
 \operatorname{RMSE}=\sqrt{\frac1N\sum_{i=1}^N(\widehat b_i-b_i)^2}
\]
衡量数值预测的误差，通常越小越好；分类的 AUC 是 ROC 曲线下面积，反映跨阈值的排序区分能力，并不等于选定阈值下的准确率。不能把某一指标的数值直接解释为另一指标的百分比。数据应用应同时考虑预测用途、验证结果与隐私影响。

\subsection{数据、模型与算法的配合}
机器翻译可以依靠人工编写词法、语法规则，也可以从大量对齐语料中学习语言之间的对应。后一种途径用数据和计算补充难以逐条编写的规则，说明复杂任务可以通过简单模型的组合或学习得到的表示来处理。

围棋程序则同时需要局面判断与行动搜索。训练数据帮助建立局面评价和落子选择的模型；蒙特卡罗树搜索通过对后续行动进行选择、扩展、模拟或评估，并回传结果，逐步把计算集中到更有希望的分支。模型提供判断，搜索组织决策过程，二者配合才形成完整系统。

这些案例表明，数据本身并不直接给出答案。还需要确定表示、模型、学习方法和计算过程，才能把记录转化为可使用的预测。

\section{机器学习与深度学习}
\subsection{人工智能与学习方法的层次}
人工智能关注智能任务的实现，机器学习则通过数据学习其中的规律。语音识别、图像识别、翻译和问答都属于人工智能的任务；完成这些任务可以采用知识规则，也可以采用学习方法。

知识库方法把事实和规则明确编码，再用逻辑规则推理。它便于表达清楚的知识，但对难以形式化描述的感知任务，逐条编写规则十分困难。机器学习从样本中学习模式，可以用于分类、预测和发现结构，但其效果仍受数据和模型限制。

图灵测试通过对话中能否区分机器与人来讨论机器智能。这是一种特定的行为评价方式；某项识别任务做得好，并不能据此证明机器具备人的全部智能。

几类方法之间具有如下关系：
\[
 \text{深度学习}\subset\text{表示学习}\subset\text{机器学习}\subset\text{人工智能}.
\]
表示学习让模型从数据中学习有效表示；深度学习进一步把简单表示逐层组合为复杂表示。并非所有机器学习都需要深层网络，也并非所有人工智能方法都通过数据训练。

\subsection{样本、特征与学习任务}
要让模型利用数据，需要先确定每个样本怎样表示，以及要预测什么。记训练数据为
\[
 D=\{(\va_i,b_i):i=1,\ldots,N\},\qquad \va_i\in\R^d.
\]
其中 $\va_i$ 是第 $i$ 个样本的特征向量，$b_i$ 是相应标签或响应。模型参数与输入特征是不同对象：特征描述样本，参数则在训练中被调整，用来建立输入与输出之间的联系。
\begin{itemize}
\item \textbf{监督学习：}利用带标签样本学习输入到输出的关系。输出为连续数值时称为回归，输出为类别时称为分类。
\item \textbf{无监督学习：}样本没有给定目标标签，任务是发现聚类、低维表示等结构。聚类得到的分组不同于已知类别下的分类。
\item \textbf{强化学习：}通过与环境互动，利用动作产生的奖励学习策略，关注长期行动效果。奖励不同于每个样本的正确类别标签。
\end{itemize}

特征选择从已有特征中选取子集；特征提取通过变换构造新的特征；表示学习让这种变换也由数据学习。例如，手工设计图像边缘特征属于特征提取，网络逐层学到边缘、局部形状和更抽象结构则体现表示学习。特征应服务于目标，维度多并不必然优于维度少。

表示方式还会改变分类的难易。考虑同心分布的两类平面点：内层点云被外层环状点云围住时，在原来的笛卡尔坐标中不能用一条直线将它们分开。若改用极坐标，取点到共同中心的距离 $r$ 作为特征，就可用一个半径阈值区分内外两类。改变的是表示及其下的分界形式，原来的样本和类别并没有改变。

\subsection{训练、损失与泛化}
训练通过调整参数减小损失，泛化则要求模型在未用于训练的数据上仍有较好表现。设 $f(\va;\vx)$ 为参数 $\vx$ 下的预测，$\ell$ 为单个样本的损失，常见训练模型为
\begin{equation}\label{eq:learning}
 \boxed{\min_{\vx\in\mathcal W}\frac1N\sum_{i=1}^N
 \ell\bigl(f(\va_i;\vx),b_i\bigr)+\lambda R(\vx).}
\end{equation}
第一项是经验损失，衡量训练数据上的拟合程度；$R(\vx)$ 是正则项，表达对参数大小、稀疏性等结构的要求；$\lambda\ge0$ 控制两者的相对权重，$\mathcal W$ 是允许的参数集合。将损失之和改为平均值时，正则参数也应相应调整，才能保持同一个优化问题。

在总体分布上希望减小的是期望损失
\[
 \E\bigl[\ell(f(\va;\vx),b)\bigr].
\]
经验损失由有限样本近似这一期望，因而训练误差小不自动保证总体误差小。

\textbf{欠拟合}指模型未能充分描述数据中的主要结构；\textbf{过拟合}指模型过多适应训练样本中的偶然细节，导致新数据上的表现较差。增加模型复杂性可能缓解欠拟合，却也可能增加过拟合风险。正则化改变选解标准，优化算法则负责求解给定的训练问题；把训练目标降得更低，并不总能改善泛化。

训练集用于估计参数，验证集用于选择模型和正则参数等设置，测试集用于评价最终方案。交叉验证把数据轮流划分为训练部分和验证部分，以比较候选设置；反复根据测试结果调参，会使测试集失去独立评价作用。对时间序列，还应保持时间先后，避免把未来信息用于过去的预测。

\subsection{神经网络的组成}
神经网络把简单的仿射变换与非线性激活逐层复合。多层感知机由输入层、若干隐藏层和输出层组成，信息沿前向连接传播。记第 $l$ 层的输出为 $\boldsymbol h^{(l)}$，则一层运算可写为
\[
 \boldsymbol h^{(0)}=\va,\qquad
 \boldsymbol z^{(l)}=W^{(l)}\boldsymbol h^{(l-1)}+\boldsymbol c^{(l)},\qquad
 \boldsymbol h^{(l)}=t_l(\boldsymbol z^{(l)}).
\]
$W^{(l)}$ 是连接权重，$\boldsymbol c^{(l)}$ 是偏置，$t_l$ 是激活函数，通常逐分量作用。所有层的权重和偏置共同构成训练参数。

常见激活函数为
\[
 \operatorname{sigmoid}(z)=\frac1{1+e^{-z}},\qquad
 \operatorname{ReLU}(z)=\max\{0,z\},\qquad
 \operatorname{Heaviside}(z)=
 \begin{cases}1,&z\ge0,\\0,&z<0.\end{cases}
\]
非线性激活使网络能够表达非线性关系。若各层都只是仿射变换，则多层复合仍然只是一次仿射变换，增加层数并未获得同样的非线性表达能力。输出层的形式还要与回归或分类任务相适应。

\subsection{不同网络处理不同结构}
图像具有局部空间结构，序列则包含前后依赖，因此网络结构应与数据结构相适应。

\textbf{卷积神经网络（CNN）}用局部连接和权重共享提取空间特征。一个卷积核在图像上滑动，对局部区域与核的对应元素相乘求和，得到特征图。若图像为 $I$、核为 $K$，一种常用计算形式为
\[
 S_{ij}=\sum_{u=1}^k\sum_{v=1}^k K_{uv}I_{i+u-1,j+v-1}.
\]
输出尺寸由核大小、步长和边界填充共同决定。与全连接相比，局部连接减少连接数，权重共享使同一种局部模式能在不同位置被识别；池化可进一步汇总局部特征。严格按信号处理术语，上式为互相关，神经网络中通常也称为卷积。

\textbf{转置卷积}使用与某个卷积线性映射对应的转置映射，常用于增加空间分辨率和生成图像。它并不是一般意义上恢复原图的逆卷积，也不能保证把丢失信息自动还原。

\textbf{循环神经网络（RNN）}通过隐藏状态保留序列上下文，例如
\[
 \boldsymbol h_t=t\bigl(W_h\boldsymbol h_{t-1}+W_a\va_t+\boldsymbol c\bigr).
\]
当前状态同时依赖当前输入与前一状态。沿时间展开后，可以把它看成参数共享的计算链，适合语音、文字等序列任务。这里的状态记忆不同于程序显式存储全部历史。

语音识别可以把观测声音看作待解码的信号，再结合语言规律选择较可能的文字。基于资料的问答则需要先限定可回答范围，检索候选片段，再判断相关性并组织回答。二者都说明：先确定任务和表示，才能选择合适的模型与学习过程。

\section{最优化问题与建模方法}
\subsection{推断与决策}
学习模型给出预测，行动选择还需要考虑资源与目标。例如，预测某条道路会拥堵以后，仍要比较可行路线的时间和成本；预测用户偏好以后，还要决定展示哪些物品。

\textbf{数据科学}主要从已知记录推断未知信息，涉及收集、清洗、表示、建模和评价。\textbf{运筹学}主要研究在给定目标和约束下怎样决策，通过数学模型比较行动方案。二者相互支持：预测为决策提供参数，优化用于估计学习模型，学习也可用于近似复杂决策中的价值或策略。

俾斯麦海战例把侦察方向与航线选择看成双方策略。一般地，设一方希望增大收益，对方希望压低这一收益。若 $a_{ij}$ 表示前者选择策略 $i$、后者选择策略 $j$ 时前者的收益，则保守决策先比较各策略的最坏收益，再取其中最大的一个：
\[
 \max_i\min_j a_{ij}.
\]
希望压低收益的一方则相应比较 $\min_j\max_i a_{ij}$。当某个位置同时是所在行的最小值、所在列的最大值时，形成纯策略鞍点；一般情况下二者未必相等。这个模型强调\textbf{在对方也会行动的条件下比较策略}，不能只挑某一种有利情形下收益最大的方案。

反潜与护航案例则从损失、搜索效果等观测出发，改进编队、搜索和资源使用。它与博弈案例共同说明：决策需要明确目标、可选行动、限制条件和比较标准。

\subsection{决策变量、目标函数与约束}
建立优化模型，就是确定可以改变的量、比较方案的标准和必须满足的条件。一般最小化问题写为
\begin{equation}\label{eq:optimization}
 \boxed{\min_{\vx\in\mathcal X} f(\vx).}
\end{equation}
\begin{definition}
$\vx$ 称为\textbf{决策变量}，$f$ 称为\textbf{目标函数}，$\mathcal X$ 称为\textbf{可行域}。属于 $\mathcal X$ 的点称为可行点；若 $\vx^*\in\mathcal X$ 且对所有 $\vx\in\mathcal X$ 有
\[
 f(\vx^*)\le f(\vx),
\]
则 $\vx^*$ 为全局最优解，$f(\vx^*)$ 为最优值。
\end{definition}
可行域常由等式与不等式组成，例如
\[
 \mathcal X=\{\vx:g_i(\vx)\le0,\ i=1,\ldots,p;
 \ h_j(\vx)=0,\ j=1,\ldots,q\}.
\]
无约束问题取 $\mathcal X=\R^n$。最大化 $f$ 可以改写为最小化 $-f$，两者的最优解相同，最优值互为相反数。

最优解不一定存在。例如 $\min_{x\in\R}e^x$ 的下确界是 $0$，却没有任何有限 $x$ 取得这个值。可行、达到最优与只有下确界，是不同的情况；目标函数也可能向下无界。

\begin{example}[投资组合]
设 $x_i$ 为投在第 $i$ 项资产的资金比例，$\mu_i$ 为预期回报率，$\Sigma$ 为回报率的协方差矩阵。以方差度量风险，在不允许卖空且要求预期回报至少为 $r_0$ 时，可写为
\[
 \begin{aligned}
 \min_{\vx}\quad&\frac12\vx^{\mathsf T}\Sigma\vx,\\
 \text{s.t.}\quad&\boldsymbol\mu^{\mathsf T}\vx\ge r_0,\qquad
 \sum_{i=1}^n x_i=1,\qquad x_i\ge0.
 \end{aligned}
\]
\end{example}
资金比例是变量，风险是目标，预算、收益和非负要求是约束。风险还可用风险价值（VaR）、条件风险价值（CVaR）等指标衡量：VaR 是给定置信水平下损失的分位数；在连续分布情形下，CVaR 可理解为损失越过这一分位点后的尾部平均损失。不同指标可能使模型具有不同的光滑性与计算性质。因此，实际目标相似，不意味着数学模型和求解算法相同。

\subsection{模型的类别与性质}
同一个实际问题可以有不同模型，而模型的性质会影响求解难度。
\begin{itemize}
\item \textbf{线性规划：}目标为线性函数，等式、不等式约束也是线性的。常数项不影响这一分类。
\item \textbf{二次规划：}目标为二次函数，约束为线性等式或不等式。投资组合的方差模型属于此类。
\item \textbf{非线性规划：}目标或约束包含非线性关系。二次规划也是非线性规划的一类。
\item \textbf{整数与组合优化：}部分或全部变量只能取整数，或表示离散选择。
\item \textbf{非光滑优化：}目标或约束含不可微结构，如绝对值在零点不可微。无导数方法则指不直接利用导数信息的方法，所处理的函数不一定不可微。
\item \textbf{半定规划：}在线性等式、不等式及矩阵半正定约束下优化线性目标。对称矩阵 $X\succeq0$ 表示对任意向量 $\boldsymbol v$，都有 $\boldsymbol v^{\mathsf T}X\boldsymbol v\ge0$。
\end{itemize}
稀疏优化与低秩优化按希望得到的解结构命名；随机优化考虑随机变量或期望目标，鲁棒优化考虑不确定条件下的性能。这些分类可以交叉，不能看成互不相容的几类。

\begin{definition}[凸性]
若对 $\vx,\boldsymbol y\in\mathcal X$ 和 $t\in[0,1]$，都有 $t\vx+(1-t)\boldsymbol y\in\mathcal X$，则 $\mathcal X$ 为凸集。若在凸集上
\[
 f(t\vx+(1-t)\boldsymbol y)\le t f(\vx)+(1-t)f(\boldsymbol y),
\]
则 $f$ 为凸函数。在凸集上最小化凸函数，称为凸优化问题。
\end{definition}
凸优化的重要性质是局部最优解也是全局最优解，但凸性本身不保证解一定存在或唯一。非凸问题可能有不同的局部最优解，算法对初始点的依赖通常更值得关注。

\subsection{目标与约束的设计}
观测通常含有误差，目标函数需要表达怎样衡量误差，以及希望解具有什么结构。设预测为 $\phi_i(\vx)$、观测为 $b_i$，残差为 $r_i(\vx)=\phi_i(\vx)-b_i$。三种常见选择为
\[
 \min_{\vx}\sum_i r_i(\vx)^2,\qquad
 \min_{\vx}\sum_i|r_i(\vx)|,\qquad
 \min_{\vx}\max_i|r_i(\vx)|.
\]
平方损失对大残差给予更大的惩罚；绝对值损失相对减弱极端残差的影响；最大残差目标直接控制最坏误差。选择哪一种，应由误差特点和应用要求决定。

成本、损失和收益也可直接作为目标。若变量是函数，目标可能是定义在函数空间上的泛函，例如能量；通过离散化，可把某些问题转成有限维优化问题。约束则表达资源、非负性、守恒关系、正交关系或其他物理要求。模型合理性不仅表现为拟合已知数据，还要看能否反映实际限制，并与已有方法作公平比较。

\subsection{等价变换}
有时模型难以直接求解，可以先改变表达形式，同时保留原问题的含义。例如
\[
 \min_{\vx} f(A\vx+\vb)
 \quad\Longleftrightarrow\quad
 \min_{\vx,\boldsymbol y} f(\boldsymbol y),\quad
 \text{s.t. }\boldsymbol y=A\vx+\vb.
\]
右边引入了新变量，但没有放宽 $\boldsymbol y$ 与 $\vx$ 的关系。同样，$\min_{\vx}\{h(\vx)+r(\vx)\}$ 可拆为
\[
 \min_{\vx,\boldsymbol y}h(\vx)+r(\boldsymbol y),
 \qquad\text{s.t. }\vx=\boldsymbol y.
\]
对不等式 $g(\vx)\le0$，引入非负变量 $s$，可等价写为 $g(\vx)+s=0,\ s\ge0$。这个 $s$ 虽称为松弛变量，但该改写\textbf{没有扩大原变量的可行域}。

引入目标值的上界变量 $t$，$\min_{\vx} f(\vx)$ 还可等价写为
\[
 \min_{\vx,t}t,\qquad\text{s.t. }f(\vx)\le t.
\]
等价变换通过消去辅助变量可回到原问题；它改变计算组织方式，而不改变原变量下的可行方案和最优目标。

\subsection{替代目标与约束松弛}
另一类处理会改变原问题。例如，用 $\ell_1$ 范数替代非零分量计数，或用核范数替代矩阵秩，保留了促进稀疏或低秩的意图，却一般不能保证解完全相同。

约束松弛则把可行域 $\mathcal X$ 扩大为 $\mathcal X_R\supseteq\mathcal X$。例如用 $0\le x\le1$ 替代 $x\in\{0,1\}$，或把等式约束放宽为单侧不等式。对同一最小化目标有
\[
 \inf_{\vx\in\mathcal X_R}f(\vx)\le
 \inf_{\vx\in\mathcal X}f(\vx).
\]
松弛得到的最优值提供下界，但解可能不满足原约束。若采用下界函数 $f_R\le f$ 替代目标，也能形成相应下界；并非所有启发式替代目标都自动具有这一性质。

对于单位球面约束，可以把 $X=\vx\vx^{\mathsf T}$ 看成提升后的变量。其精确矩阵描述为
\[
 X\succeq0,\qquad\operatorname{tr}(X)=1,\qquad\rank X=1.
\]
这里 $\operatorname{tr}(X)$ 是对角元素之和，等于 $\norm{\vx}_2^2$。去掉秩为 $1$ 的约束，就得到更大的集合 $X\succeq0,\ \operatorname{tr}(X)=1$。若结果仍为秩 $1$，可以分解回单位向量；若秩大于 $1$，则不能直接写成单个 $\vx\vx^{\mathsf T}$。在具体模型中，还须把目标与其他约束相应表达在提升变量下，才构成完整松弛。

\section{回归、似然与正则化}
\subsection{最小二乘与回归}
回归用输入特征描述连续输出，最小二乘则用残差平方和衡量拟合误差。一般回归模型为
\[
 b_i=f(\va_i;\vx)+\varepsilon_i,
\]
其中输入 $\va_i$ 与输出 $b_i$ 已观测，参数 $\vx$ 待估计，误差 $\varepsilon_i$ 未知。把函数写成含有限个参数的形式，就把选择函数的问题转成选择参数的问题。

在线性模型中，令 $A$ 的第 $i$ 行为 $\va_i^{\mathsf T}$，$\vb=(b_1,\ldots,b_N)^{\mathsf T}$，则
\[
 \vb=A\vx+\boldsymbol\varepsilon,\qquad
 \boxed{\min_{\vx}\frac12\norm{A\vx-\vb}_2^2.}
\]
这里“线性”指对待估参数线性。最小二乘最优解满足正规方程
\[
 A^{\mathsf T}A\vx=A^{\mathsf T}\vb.
\]
若 $A$ 列满秩，则解唯一；否则参数可能不唯一。即使原方程 $A\vx=\vb$ 无解，仍可寻找残差平方和最小的近似解。非线性预测函数也可以使用平方损失，此时得到非线性最小二乘问题。

\subsection{最大似然估计}
最大似然估计选择使已观测样本更可能出现的参数。设样本 $z_1,\ldots,z_N$ 独立同分布，概率质量或密度为 $p(z;\vx)$，似然函数为
\[
 L(\vx)=\prod_{i=1}^N p(z_i;\vx).
\]
观测样本固定，似然作为参数 $\vx$ 的函数。对连续变量使用的是密度，并非每个精确观测点有正概率；似然也不是“参数本身的概率”。

由于对数严格递增，最大化 $L$ 与最大化 $\log L$ 等价，常进一步改写为最小化负对数似然：
\[
 \max_{\vx}L(\vx)
 \quad\Longleftrightarrow\quad
 \min_{\vx}-\sum_{i=1}^N\log p(z_i;\vx).
\]
乘积变成求和，既便于分析，也便于数值求解；但最大似然解仍可能不存在或不唯一。

若线性回归误差独立且服从共同方差的零均值高斯分布，则
\[
 p(b_i\mid\va_i;\vx)
 =\frac1{\sqrt{2\pi\sigma^2}}
 \exp\!\left(-\frac{(b_i-\va_i^{\mathsf T}\vx)^2}{2\sigma^2}\right).
\]
对给定的共同方差 $\sigma^2>0$，负对数似然等于与 $\vx$ 无关的常数加上
\[
 \frac1{2\sigma^2}\sum_{i=1}^N(b_i-\va_i^{\mathsf T}\vx)^2.
\]
因此参数 $\vx$ 的最大似然估计就是最小二乘解，求这个解不需要事先知道方差的具体数值。若误差独立服从共同尺度的零中心 Laplace 分布，相应最大似然则最小化绝对残差和。\textbf{回归任务、误差假设与损失函数相互关联，但不是同一个概念。}

\subsection{岭回归与LASSO}
仅减小训练残差，可能得到不稳定或过于复杂的解，因此要同时控制参数。对 $\vx=(x_1,\ldots,x_n)^{\mathsf T}$，记
\[
 \norm{\vx}_0=\#\{j:x_j\ne0\},\qquad
 \norm{\vx}_1=\sum_j|x_j|,\qquad
 \norm{\vx}_2=\left(\sum_jx_j^2\right)^{1/2}.
\]
$\norm{\vx}_0$ 只是非零分量的计数，通常称为 $\ell_0$“范数”，严格说并不是范数。

\textbf{岭回归}采用 $\ell_2$ 范数平方正则项，是 Tikhonov 正则化的一种形式：
\begin{equation}\label{eq:ridge}
 \boxed{\min_{\vx}\frac12\norm{A\vx-\vb}_2^2+\lambda\norm{\vx}_2^2,
 \qquad\lambda>0.}
\end{equation}
其最优解满足
\[
 (A^{\mathsf T}A+2\lambda I)\vx=A^{\mathsf T}\vb.
\]
由于 $A^{\mathsf T}A+2\lambda I$ 正定，模型有唯一解。正则项抑制过大的系数，改善病态问题中的稳定性；一般不会使大量系数恰好为零。这里所有分量均被正则化；若另设截距而不惩罚截距，应按相应参数分块处理。

\textbf{LASSO}采用 $\ell_1$ 正则项：
\begin{equation}\label{eq:lasso}
 \boxed{\min_{\vx}\frac12\norm{A\vx-\vb}_2^2+\lambda\norm{\vx}_1,
 \qquad\lambda>0.}
\end{equation}
它在拟合数据的同时促进部分系数为零，从而进行变量选择。在二维中，$\ell_1$ 单位球是顶点位于坐标轴上的菱形，$\ell_2$ 单位球则是圆。在相应的约束形式下，前者的顶点使损失等高线容易在某个坐标为零的位置接触约束边界，从而产生稀疏解。这也与 $\ell_1$ 范数在含零坐标的位置可能不可微相联系。两种惩罚的作用不同：$\ell_1$ 促进稀疏，$\ell_2$ 惩罚则倾向于连续缩小系数。

正则参数越大，目标中对结构或参数大小的要求相对越强；过大可能造成欠拟合，过小则可能抑制不足。参数应结合验证表现选择，而不是只追求训练误差最小。$\ell_1$ 促进稀疏，不保证任意数据、任意参数下都能恢复真实非零位置。

\subsection{惩罚形式与约束形式}
正则化还可写成约束。以下两类模型分别限制参数大小和限制残差：
\[
 \min_{\vx}\frac12\norm{A\vx-\vb}_2^2,
 \quad\text{s.t. }\norm{\vx}_1\le\tau;
\]
\[
 \min_{\vx}\norm{\vx}_1,
 \quad\text{s.t. }\norm{A\vx-\vb}_2\le\delta.
\]
$\tau$ 表示参数预算，$\delta$ 表示可接受的拟合误差。岭回归也有 $\norm{\vx}_2\le\tau$ 的约束形式。惩罚式与约束式都表达拟合与结构之间的权衡；在相应条件和参数对应关系下可以共享最优解，但不能把任意 $\lambda,\tau,\delta$ 直接视为等价。

还可以改变损失本身，例如
\[
 \min_{\vx}\norm{A\vx-\vb}_2+\lambda\norm{\vx}_1,
 \qquad
 \min_{\vx}\norm{A\vx-\vb}_1+\lambda\norm{\vx}_1.
\]
前者使用未平方的 $\ell_2$ 残差，后者使用绝对残差；二者与标准 LASSO 的损失不同，所对应的误差假设和计算性质也可能不同。

\subsection{结构化正则化}
若特征之间存在分组或变换结构，正则项也应反映这种先验。将指标划分为 $G$ 个不相交的组 $I_1,\ldots,I_G$，记 $n_g=|I_g|$，分组 LASSO 为
\[
 \min_{\vx}\frac12\norm{A\vx-\vb}_2^2
 +\lambda\sum_{g=1}^G\sqrt{n_g}\norm{\vx_{I_g}}_2.
\]
它对各组内部先取 $\ell_2$ 范数，再将组范数相加，促进整组变量为零。组被选中不代表组内每个分量都一定非零；该模型主要控制组的选择，不特别追求组内稀疏。当各组只有一个变量时，退化为普通 LASSO。

同时要求组与组内的稀疏性，可以使用\textbf{稀疏分组 LASSO}：
\[
 \min_{\vx}\frac12\norm{A\vx-\vb}_2^2
 +\lambda_1\sum_{g=1}^G\sqrt{n_g}\norm{\vx_{I_g}}_2
 +\lambda_2\norm{\vx}_1.
\]
若稀疏性出现在某个变换域，则采用
\[
 \min_{\vx}\frac12\norm{A\vx-\vb}_2^2+\lambda\norm{F\vx}_1,
 \qquad F\in\R^{p\times n}.
\]
例如 $F$ 取相邻分量差分算子时，$\norm{F\vx}_1=\sum_{j=2}^n|x_j-x_{j-1}|$，使相邻差异倾向于稀疏，适于描述分段平稳结构。\textbf{融合 LASSO}进一步同时惩罚系数及其差分：
\[
 \min_{\vx}\frac12\norm{A\vx-\vb}_2^2
 +\lambda_1\norm{\vx}_1
 +\lambda_2\sum_{j=2}^n|x_j-x_{j-1}|.
\]
信号本身稀疏、分组稀疏与变换后稀疏是不同假设，不能不加区分地使用同一个正则项。

\section{稀疏优化与低秩恢复}
\subsection{欠定方程中的稀疏解}
当未知数多于独立方程时，观测不足以唯一确定解；稀疏性为选解提供额外结构。设 $A\in\R^{m\times n}$，$m<n$。若 $A\vx=\vb$ 有解，则由于 $\rank A<n$，解不唯一；若不满足相容条件，也可能无解。

在已知真实信号只有少数非零分量的情况下，可考虑
\[
 \min_{\vx}\norm{\vx}_0,\qquad\text{s.t. }A\vx=\vb.
\]
该模型直接寻找非零分量最少的解，但一般计算困难。用凸的 $\ell_1$ 范数替代，得到\textbf{基追踪}模型
\[
 \boxed{\min_{\vx}\norm{\vx}_1,\qquad\text{s.t. }A\vx=\vb.}
\]
在测量矩阵、信号稀疏程度等满足适当条件时，基追踪能够恢复真实稀疏信号。压缩感知正是利用这些结构条件，从较少测量中恢复较高维信号；它并不是任意少量观测都能还原任意信息。

基追踪还可等价写成线性规划：引入 $t_j\ge0$，用 $-t_j\le x_j\le t_j$ 限制各分量，再最小化 $\sum_j t_j$，并保留 $A\vx=\vb$。不过，等价于线性规划并不意味着通用线性规划算法总适合大规模计算。当 $A$ 稠密，或只能通过算子计算 $A\vx$ 与 $A^{\mathsf T}\boldsymbol y$ 时，还要考虑存储、矩阵运算及算法能否利用这种结构。

若改为最小化 $\norm{\vx}_2$，得到的是原点到仿射解集的最短距离点，一般不追求稀疏。若 $A$ 行满秩，最小 $\ell_2$ 范数解可写为
\[
 \vx=A^{\mathsf T}(AA^{\mathsf T})^{-1}\vb.
\]
这与基追踪选择解的原则不同。带噪观测则常用误差约束或 LASSO 的惩罚形式，避免强迫模型精确拟合噪声。

\subsection{低秩矩阵补全}\label{sec:lowrank}
用户偏好可能由少数共同因素解释，因此可以用低秩结构连接已知评分与未知评分。设 $M\in\R^{m\times n}$ 的已观测位置为
\[
 \Omega=\{(i,j):M_{ij}\text{ 已知}\},
\]
用 $X$ 表示待求的完整评分矩阵，理想化模型为
\[
 \boxed{\min_X\rank X,\qquad
 \text{s.t. }X_{ij}=M_{ij},\quad(i,j)\in\Omega.}
\]
约束保留已知评分，目标用少数潜在因素解释整个矩阵。\textbf{观测稀疏}指已知位置少，\textbf{矩阵低秩}指行列之间存在低维线性结构；两者不是同一性质。低秩是假设，也不单独保证补全唯一。

设 $X$ 的奇异值为 $\sigma_1,\ldots,\sigma_{\min(m,n)}$，则
\[
 \rank X=\#\{j:\sigma_j>0\},\qquad
 \norm{X}_*=\sum_j\sigma_j.
\]
$\norm{X}_*$ 称为核范数。矩阵秩应由非零奇异值计数，一般不能直接用非零特征值个数替代。用核范数替代秩，得到凸模型
\[
 \boxed{\min_X\norm{X}_*,\qquad
 \text{s.t. }X_{ij}=M_{ij},\quad(i,j)\in\Omega.}
\]
核范数促进低秩，能否精确恢复还依赖观测数量、位置和真实矩阵的结构，不能认为替代后总与原问题等价。

观测含噪声时，可采用
\[
 \min_X\frac12\sum_{(i,j)\in\Omega}(X_{ij}-M_{ij})^2
 +\lambda\norm{X}_*.
\]
残差只在已观测位置计算；若把缺失位置当作零分加入损失，就改变了原问题。

\subsection{低秩分解与两类结构的联系}
低秩矩阵还可以表示为两个较窄矩阵的乘积，从而减少待处理的参数。取
\[
 X=UV^{\mathsf T},\qquad U\in\R^{m\times r},\quad V\in\R^{n\times r},
 \qquad r\ll\min(m,n),
\]
则 $\rank X\le r$，不一定恰好等于 $r$。可直接求解
\[
 \min_{U,V}\frac12\sum_{(i,j)\in\Omega}
 \bigl((UV^{\mathsf T})_{ij}-M_{ij}\bigr)^2
 +\alpha\norm{U}_F^2+\beta\norm{V}_F^2,
\]
其中 $\norm{U}_F^2$ 表示全部元素平方之和，$\alpha,\beta>0$。参数量由 $mn$ 减为约 $(m+n)r$，但因 $U,V$ 相乘，联合问题通常非凸。

正则项可以抑制因子尺度过大，但不会消除一切表示不唯一性。例如对正交矩阵 $Q$，$(UQ)(VQ)^{\mathsf T}=UV^{\mathsf T}$，两因子的 Frobenius 范数也不变。

稀疏与低秩具有平行关系：向量的 $\ell_0$ 计数对应矩阵的非零奇异值计数，向量的 $\ell_1$ 范数对应奇异值之和，即核范数。两者都借助较简单的结构，从不充分或含噪观测中选择较合理的解。

\section{逻辑回归与支持向量机}
\subsection{二分类与逻辑回归}
分类要判断样本所属类别，逻辑回归通过概率描述二分类结果。设标签 $b_i\in\{-1,1\}$，令
\[
 s(\va)=\vw^{\mathsf T}\va+c,\qquad \sigma(z)=\frac1{1+e^{-z}}.
\]
逻辑回归规定
\[
 \Pr(b=1\mid\va)=\sigma(s(\va)),\qquad
 \Pr(b=-1\mid\va)=\sigma(-s(\va)).
\]
即 $\Pr(b\mid\va)=\sigma(b\,s(\va))$。参数为 $\vw,c$；也可把常数 $1$ 加入特征向量，将截距吸收到参数中。

在样本独立的条件下，最大化标签的条件似然，等价于最小化
\begin{equation}\label{eq:logistic}
 \boxed{\sum_{i=1}^N\log\bigl(1+\exp(-b_i s(\va_i))\bigr).}
\end{equation}
这是逻辑损失，也可加 $\ell_2$ 或 $\ell_1$ 正则项。若使用 $0,1$ 标签 $y_i$ 和预测概率 $p_i$，同一负对数似然为
\[
 -\sum_i\bigl[y_i\log p_i+(1-y_i)\log(1-p_i)\bigr],
\]
即二分类交叉熵损失，这里通常用自然对数。

“回归”是方法名称，逻辑回归通常用于分类；它预测类别概率，线性回归则通常预测连续响应。在线性评分下，以 $p=1/2$ 为分类阈值，对应决策边界 $s(\va)=0$。无正则时，线性可分数据上的逻辑损失可能只有下确界而没有有限参数最优解，不能默认任何训练问题都能取得最小值。

\subsection{支持向量机与最大间隔}
仅把两类训练点分开，并不能确定唯一的分类平面；支持向量机进一步选择间隔较大的平面。考虑 $\vw\ne0$ 的超平面
\[
 \vw^{\mathsf T}\va+c=0.
\]
点 $\va_i$ 到它的距离为 $|\vw^{\mathsf T}\va_i+c|/\norm{\vw}_2$。若分类正确，则 $b_i(\vw^{\mathsf T}\va_i+c)>0$；带标签的几何间隔为
\[
 \frac{b_i(\vw^{\mathsf T}\va_i+c)}{\norm{\vw}_2}.
\]
最大间隔方法使所有样本中最小的这个距离尽量大，让分类边界尽量远离最靠近它的训练点。

同时用一个正数缩放 $\vw,c$，不会改变分类平面与几何距离。因此在线性可分时，可以把最小函数间隔规范为 $1$，要求
\[
 b_i(\vw^{\mathsf T}\va_i+c)\ge1.
\]
此时最近点到分界面的距离为 $1/\norm{\vw}_2$，两条支撑超平面之间的总宽度为 $2/\norm{\vw}_2$。最大化间隔等价于\textbf{硬间隔支持向量机}：
\[
 \boxed{\begin{aligned}
 \min_{\vw,c}\quad&\frac12\norm{\vw}_2^2,\\
 \text{s.t.}\quad&b_i(\vw^{\mathsf T}\va_i+c)\ge1,
 \quad i=1,\ldots,N.
 \end{aligned}}
\]
硬间隔模型中，位于间隔边界上的训练点是支持向量，约束了分界面的位置。最大间隔使分类对一定程度的位置扰动更有余地，但不能据此保证对任意新数据或噪声都正确。

\subsection{软间隔与分类方法比较}
实际数据未必线性可分，因此允许个别样本违反间隔要求，并对违反程度施加惩罚。引入 $\xi_i\ge0$，得到\textbf{软间隔支持向量机}：
\[
 \boxed{\begin{aligned}
 \min_{\vw,c,\boldsymbol\xi}\quad&\frac12\norm{\vw}_2^2+C\sum_{i=1}^N\xi_i,\\
 \text{s.t.}\quad&b_i(\vw^{\mathsf T}\va_i+c)\ge1-\xi_i,\qquad\xi_i\ge0,
 \end{aligned}}
 \qquad C>0.
\]
固定 $\vw,c$ 后，最优松弛量为
\[
 \xi_i=\max\{1-b_i s(\va_i),0\}.
\]
因此可根据带标签评分解释违反程度：
\begin{itemize}
\item $b_i s(\va_i)\ge1$：分类正确且满足间隔，$\xi_i=0$；
\item $0<b_i s(\va_i)<1$：分类正确但进入间隔内部，$0<\xi_i<1$；
\item $b_i s(\va_i)=0$：位于分界面，$\xi_i=1$；
\item $b_i s(\va_i)<0$：被误分类，$\xi_i>1$。
\end{itemize}
所以 $\sum_i\xi_i$ 惩罚的是\textbf{间隔违反程度}，不是误分类点的个数；$\xi_i>0$ 也不一定表示误分类。

消去松弛变量，得到等价的合页损失形式：
\[
 \min_{\vw,c}\frac12\norm{\vw}_2^2
 +C\sum_i\max\{1-b_i(\vw^{\mathsf T}\va_i+c),0\}.
\]
$C$ 越大，违反间隔的代价相对越高；$C$ 较小则更容许违反，以换取较小的权重范数。若希望进行特征选择，还可将平方 $\ell_2$ 正则替换为 $\ell_1$ 正则，得到
\[
 \min_{\vw,c}\norm{\vw}_1
 +C\sum_i\max\{1-b_i(\vw^{\mathsf T}\va_i+c),0\}.
\]
这一模型促进特征系数稀疏，但已改变了原来基于欧氏距离的最大间隔模型。

逻辑回归与线性 SVM 都可得到线性决策边界，但依据不同：逻辑回归从条件概率与最大似然出发，使用光滑的逻辑损失；SVM 从最大间隔出发，使用合页损失。前者直接给出模型概率，后者的原始评分不直接等于概率。两者都可加入结构正则项，不能单凭名称判断哪一种在所有任务上更好。

\section{优化算法的基本思想}
\subsection{从模型到迭代}
模型规定求什么，算法规定怎样逐步计算。目标和约束确定后，仍须选择初始点、更新规则、步长与停止准则。一般迭代形式为
\[
 \vx^{k+1}=\vx^k+\alpha_k\boldsymbol d^k,
\]
其中 $\boldsymbol d^k$ 是搜索方向，$\alpha_k>0$ 是步长。方向表达往哪里走，步长决定走多远。步长太大可能越过有利区域甚至发散，太小则可能进展缓慢。

算法应与模型结构相适应。可微目标能够使用梯度信息，含 $\ell_1$ 范数或合页损失的目标则不能在所有点直接套用普通梯度；整数约束也不能仅靠连续的小幅移动处理。即使数学模型正确，计算过程和停止判断仍需分别检验。

\subsection{梯度与随机梯度}
对可微函数，梯度由各个偏导数组成：
\[
 \nabla f(\vx)=\left(\frac{\partial f}{\partial x_1},\ldots,
 \frac{\partial f}{\partial x_n}\right)^{\mathsf T}.
\]
当梯度非零时，沿负梯度方向取充分小的步长可使目标下降。\textbf{梯度下降}采用
\[
 \boxed{\vx^{k+1}=\vx^k-\alpha_k\nabla f(\vx^k).}
\]
其思想是利用当前位置的一阶变化信息逐步改善目标；是否收敛以及收敛到什么点，还取决于目标性质和步长规则。

大规模学习的目标常为 $f(\vx)=N^{-1}\sum_{i=1}^N f_i(\vx)$。全梯度每次使用全部样本，计算较准确但代价较大。\textbf{随机梯度下降（SGD）}每次随机抽取样本，用单个或少量样本的梯度近似全梯度，例如
\[
 \boldsymbol g^k=\frac1{|B_k|}\sum_{i\in B_k}\nabla f_i(\vx^k),\qquad
 \vx^{k+1}=\vx^k-\alpha_k\boldsymbol g^k.
\]
当 $|B_k|=1$ 时使用单个样本，取若干样本则为小批量方法。适当的均匀抽样下，$\boldsymbol g^k$ 是全梯度的无偏估计；但一次更新可能使完整目标升高，不能要求每一步都单调下降。

随机方法以带噪声的方向换取较低的单步代价，适合样本量很大的问题。较大的批量通常减少梯度噪声，却增加单步计算量；总体效率还受并行计算、步长与目标结构影响。

在神经网络中，\textbf{前向传播}计算预测和损失，\textbf{反向传播}利用链式法则高效计算各层参数的梯度，优化器再用梯度更新参数。反向传播负责求导，SGD、Adam 等负责更新，不能把二者看作同一个环节。

\subsection{动量、自适应步长与其他方法}
大规模训练中，常在随机梯度基础上积累方向信息或调整各分量的步长。\textbf{动量方法}积累近期方向，减少更新的反复摆动；Nesterov 类方法进一步利用预测位置的信息。\textbf{自适应方法}根据历史梯度大小，对不同坐标采用不同尺度的更新。
\begin{center}
\begin{tabular}{p{0.36\textwidth}p{0.53\textwidth}}
\toprule
方法 & 基本认识\\
\midrule
SGD、ASGD & 随机梯度更新；ASGD 还对迭代结果作平均。\\
Momentum、Nesterov & 利用方向积累或预测位置改善更新。\\
Adagrad & 累积各坐标梯度平方，调整相应学习率。\\
RMSprop、Adadelta & 利用近期平方梯度等尺度信息，避免只按全部历史累积。\\
Adam、Adamax & 结合方向平均与尺度调整；Adamax 使用不同的尺度估计。\\
\bottomrule
\end{tabular}
\end{center}
这些方法改变的是训练目标的求解过程，不会自动修正数据偏差，也不普遍保证找到非凸目标的全局最优解。其他算法名称如 FTRL、ProximalAdagrad、YellowFin，对应不同的结构处理或参数调整思路；理解它们属于优化器及其变种即可，不应把名称列表当作算法推导。

\textbf{牛顿法}利用二阶变化信息。对二阶可微的无约束目标，在当前点解线性方程
\[
 \nabla^2 f(\vx^k)\boldsymbol d^k=-\nabla f(\vx^k),
\]
其中 $\nabla^2 f$ 是由二阶偏导数组成的 Hessian 矩阵。求出方向后，再沿该方向更新，相当于利用局部二次模型选取方向；在适当条件且接近非退化极小点时可以很快，但单步代价较高。非凸问题的 Hessian 矩阵可能不是正定的，原始牛顿方向未必是下降方向，通常还需相应的修正或步长控制。

\textbf{原始--对偶方法}同时处理原决策变量与约束对应的乘子。乘子是反映约束作用的辅助变量，也称对偶变量；算法结合约束满足情况调整这些变量，再用它们影响原变量的更新，从而协调目标改善与约束满足。约束因而不只是更新后才检查的条件，还能参与求解过程。具体更新与收敛条件依问题结构而定。

\section{最优性、收敛与停止准则}
\subsection{全局最优与局部最优}
局部最优只比较邻近可行点，全局最优则比较全部可行点。
\begin{definition}
设 $\overline{\vx}\in\mathcal X$。若存在 $\delta>0$，使对所有满足 $\norm{\vx-\overline{\vx}}_2<\delta$ 的可行点都有
\[
 f(\overline{\vx})\le f(\vx),
\]
则称 $\overline{\vx}$ 为局部极小解。若对其中一切 $\vx\ne\overline{\vx}$ 都为严格不等式，则称为严格局部极小解。
\end{definition}
全局极小解一定是局部极小解，反之一般不成立。对凸优化问题，若存在更好的可行点，沿它与局部极小点之间的线段就能在任意近处找到更好的点，与局部极小矛盾，故局部极小必为全局极小。

在可微无约束问题的内部极小点，梯度为零；但\textbf{驻点不一定是极小点}，也可能是极大点或鞍点。例如 $f(x)=x^3$ 在 $0$ 处导数为零，却不是极小点。对带约束问题，最优点的普通梯度未必为零，例如 $\min_{x\ge0}x$ 的最优点为 $0$，梯度始终为 $1$。

\subsection{收敛对象与初始点}
评价迭代算法，既要说明什么量趋于什么，也要说明对初始点有什么要求。
\begin{itemize}
\item \textbf{点列收敛：}$\norm{\vx^k-\vx^*}\to0$，迭代点趋向某个解。
\item \textbf{函数值收敛：}$f(\vx^k)\to f^*$，目标值趋向所要求的极小值。
\item \textbf{最优性度量收敛：}例如可微无约束问题中 $\norm{\nabla f(\vx^k)}\to0$，或约束问题中适当的最优性条件违反度趋于零。
\end{itemize}
若目标连续，点列收敛到 $\vx^*$ 可推出函数值收敛到 $f(\vx^*)$；反之一般不成立。函数在多个点有同一个最优值时，点列可能在这些点之间移动。梯度范数趋于零只说明一阶残差趋于零，并不保证迭代点到真实驻点的距离趋于零。

聚点是某个子列的极限。若梯度范数趋于零，点列存在聚点且梯度连续，则该聚点为驻点，但仍不保证它是极小点或全局最优解。“每个聚点都满足某种最优性条件”也弱于“整个点列收敛到一个最优解”。若迭代过程允许不可行点，还须检查约束违反度能否随迭代减小。

算法的\textbf{局部收敛}通常要求初始点足够接近目标解，\textbf{全局收敛}则在规定的初始点范围内不需要这种邻近要求。这里的“全局”描述初始点条件；\textbf{收敛到全局最优解}描述极限解的性质，两者不能混淆。一个非凸算法即使从任意允许初始点都趋向驻点，也不能因此称其必能求得全局最优解。

\subsection{收敛速度}
在迭代确实收敛的前提下，还可以比较误差衰减的速度。设 $e_k=\norm{\vx^k-\vx^*}>0$ 且 $e_k\to0$，以下关系均考察充分大的 $k$。
\begin{itemize}
\item \textbf{Q-线性收敛：}存在 $a\in(0,1)$，使 $e_{k+1}\le a e_k$。误差每一步至少按固定比例缩小。
\item \textbf{Q-超线性收敛：}$e_{k+1}/e_k\to0$，误差缩小比例趋于零。
\item \textbf{Q-二次收敛：}存在 $M>0$，使 $e_{k+1}\le M e_k^2$；当误差很小时，下一步误差与其平方同阶受控。
\item \textbf{Q-次线性收敛：}在这里按课件的比值表述，$e_{k+1}/e_k\to1$，没有固定的小于 $1$ 的渐近收缩比例。
\end{itemize}
例如 $e_k=2^{-k}$ 为线性收敛，$e_k=2^{-2^k}$ 为二次收敛且超线性收敛，$e_k=1/k$ 为次线性收敛。这里“线性”描述误差缩小规律，不表示误差等于 $ak+b$。

若存在 Q-线性趋于零的非负序列 $t_k$，使 $e_k\le t_k$，则称为\textbf{R-线性收敛}。它用一个收敛较快的上界控制误差，允许原误差序列出现一定波动；R-超线性和 R-二次也可类似理解。

\textbf{迭代复杂度}则研究为保证达到给定精度，需要多少次迭代。例如已有保证
\[
 f(\vx^k)-f^*\le\frac{C_0}{\sqrt{k}},\qquad C_0>0,
\]
为保证误差不超过 $\varepsilon$，取 $k\ge C_0^2/\varepsilon^2$ 即可，因而得到 $O(\varepsilon^{-2})$ 的迭代次数上界。这是充分的次数保证，并非断言每次实际计算都至少需要这么多步。总运行时间还取决于单步计算、内存和通信成本，不能只比较收敛阶。

\subsection{有限步停止}
实际计算必须在有限步结束，因此需要用可计算的量判断是否已经达到所需精度。若最优值 $f^*$ 已知，可参考归一化的目标差
\[
 \frac{|f(\vx^k)-f^*|}{\max\{1,|f^*|\}}\le\varepsilon.
\]
但实际中 $f^*$ 往往未知，因此不能把它作为通用的可执行停止条件。对光滑无约束问题，梯度范数较小表示近似满足一阶驻点条件，并不直接表示到真实驻点的距离较小；对约束问题，还需检验可行性，例如
\[
 \max\{g_i(\vx^k),0\}\le\varepsilon_g,\qquad
 |h_j(\vx^k)|\le\varepsilon_h,
\]
并结合适合该问题的最优性条件，不能只检查普通梯度。

常见的停机指标还有
\[
 \frac{\norm{\vx^{k+1}-\vx^k}}{\max\{1,\norm{\vx^k}\}}\le\varepsilon_x,
 \qquad
 \frac{|f(\vx^{k+1})-f(\vx^k)|}{\max\{1,|f(\vx^k)|\}}\le\varepsilon_f.
\]
它们说明迭代或目标变化已经很小，也可能只是步长过小或算法停滞。\textbf{停止计算、满足近似最优性条件和找到全局最优解，是不同层次的结论。}最大迭代次数、时间限制等资源条件同样只能说明计算何时结束。

\section{主要概念与模型的联系}
\subsection{从数据到决策的知识结构}
数据、学习模型和优化算法分别承担不同任务，又通过同一个问题相互连接。
\begin{center}
\begin{tabular}{p{0.18\textwidth}p{0.34\textwidth}p{0.34\textwidth}}
\toprule
层次 & 核心作用 & 与下一层的联系\\
\midrule
数据与信息 & 记录现象，减少有关目标的不确定性 & 表示与质量决定模型能利用什么。\\
统计与学习 & 从样本中估计规律，预测未知结果 & 通过损失和先验把估计写为优化。\\
优化建模 & 确定变量、目标、约束与结构 & 模型性质决定可采用的算法。\\
迭代求解 & 用有限计算得到近似解 & 检验可行性、最优性和计算代价。\\
验证与决策 & 比较新数据上的效果，选择行动 & 实际反馈用于改进数据与模型。\\
\bottomrule
\end{tabular}
\end{center}
预测能力、数学求解与实际决策效果需要分别评价。例如，一个推荐模型即使准确求解了训练目标，也可能由于样本偏差而推荐不佳；一个训练误差略大的正则化模型，反而可能在新用户上更稳定。

\subsection{模型之间的比较}
比较模型时，需要同时说明输出、损失、结构假设与求解性质。
\begin{center}
\small
\begin{tabular}{p{0.18\textwidth}p{0.30\textwidth}p{0.38\textwidth}}
\toprule
模型 & 主要任务与损失 & 结构要求或主要区别\\
\midrule
最小二乘 & 连续响应；残差平方和 & 控制整体平方误差，不自动进行特征选择。\\
岭回归 & 平方损失加 $\ell_2$ 平方 & 缩小系数、改善稳定性，一般不产生大量精确零。\\
LASSO & 平方损失加 $\ell_1$ & 促进系数稀疏，涉及非光滑目标。\\
逻辑回归 & 二分类；负对数似然 & 用概率解释输出，线性评分下边界为超平面。\\
软间隔 SVM & 二分类；合页损失 & 权衡权重范数与间隔违反程度。\\
核范数补全 & 预测矩阵缺失元素 & 以奇异值之和促进低秩，仅在观测位置拟合。\\
低秩因子模型 & 矩阵拟合与因子正则 & 以较少参数表示矩阵，联合优化通常非凸。\\
\bottomrule
\end{tabular}
\end{center}

若干相近概念应保持各自含义：数据量不等于有效信息量，相关关系不等于因果关系；观测稀疏不等于矩阵低秩，参数缩小不等于特征筛选；模型是比较方案的数学表达，算法是计算方案；优化目标收敛不等于新数据预测正确，停机不等于最优。

完整说明一种方法，应使问题、原理、表达、案例与适用条件相互对应。对于大数据应用，说明记录如何形成信息、信息如何支持判断；对于优化模型，说明每个变量、目标项和约束的实际意义；对于算法，说明所用信息、更新思想与计算特点，再区分可证明的保证和实际停止条件。这样才能把概念、公式和案例连成完整解释。

从记录中取得信息，用模型形成预测，再结合目标与约束选择行动，最后通过新的观测检验结果，数据分析与最优化就在这一过程中相互配合。

\smallskip
\noindent\textbf{答题安排：}6 道论述题共 100 分钟，可以先用约 5 分钟通读题目、按分值分配时间，保留约 5 分钟检查。若各题分值相近，正文作答可按每题约 15 分钟安排；具体随分值和内容调整。作答以切合题意的定义与关系为起点，展开机制，选用相应案例或公式，最后交代成立条件，避免只罗列术语。

\clearpage
\appendix
\section{课件内容索引}
以下均为 PDF 页码。正文按知识关系重组，重复内容合并，历史数据与软件链接不作为原理的记忆要求。
\begin{longtable}{p{0.25\textwidth}p{0.15\textwidth}p{0.47\textwidth}}
\toprule
课件 & 页码 & 对应内容\\
\midrule
\endhead
大数据案例介绍 & 8--24 & 数据、信息、知识与 4V；数据与模型。\\
 & 25--39 & 推荐、流感、药物、广告与行为记录。\\
 & 40--44 & 样本量、代表性与数据驱动。\\
 & 45--54 & 博弈、运筹学及与数据科学的联系。\\
 & 55--60 & 数据价值、行为推荐、社会网络、属性预测与隐私。\\
\midrule
大数据思维 & 3--13 & 大数据特点、多维性、翻译与交通案例。\\
 & 14--18 & 确定性、不确定性、信息熵、编码、最大熵与交叉熵。\\
 & 19--26 & 异常检测、机器智能、语音与问答。\\
 & 27--33 & 知识库、机器学习、表示学习与深度学习。\\
\midrule
最优化：简介 & 3--9 & 优化的一般形式、分类与投资组合。\\
 & 11--17 & 稀疏优化、基追踪与 LASSO。\\
 & 19--23 & 低秩恢复、核范数与因子分解。\\
 & 25--35 & 学习模型、MLP、CNN、转置卷积、RNN 与优化器。\\
 & 37--45 & 最优解、收敛、速度、复杂度与停止准则。\\
\midrule
优化建模 & 2--6 & 建模流程、误差度量与正则化。\\
 & 7--10 & 最大似然、代价收益及替代目标。\\
 & 11--14 & 约束、等价变换与松弛。\\
 & 15--24 & 一般学习目标、回归、岭、LASSO 及结构变形。\\
 & 25--27 & 逻辑回归、期望风险与样本平均。\\
 & 28--31 & 硬间隔、软间隔支持向量机及合页损失。\\
\bottomrule
\end{longtable}
\end{document}
